% A page whose column holds more material than `\@colht`.
%
% `\@makecol` packs every column `\vbox to\@colht`, so a page carrying more
% than fits is set at a negative glue set: each shrinkable skip gives up its
% share and everything below the first one moves up. The float column builder
% was laying its pages out at their natural size instead, which left this page
% low by the whole shrink and 4.81pt overfull.
%
% The shrink here is `\@maketitle`'s `center` (`\@endparenv` contributes
% `\topsep` + `\partopsep` = 12pt plus 4pt minus 6pt at 11pt) plus every
% `\section`'s `3.5ex plus 1ex minus .2ex` before-skip; the stretch that makes
% stopping a line earlier look cheap is the run of one-line paragraphs, each
% preceded by `\parskip 0pt plus 1pt`. The `table` float (caption only) routes
% the document through the float column builder.
%
% Reduced from `fixtures/real-world/lmodern-report`, whose page 1 has the same
% shape: 652.76431pt of material in 650.43001pt, which `\showoutput` reports as
% `glue set - 0.26439`. No `\tableofcontents`: its page numbers need a third
% pass to settle here, which would measure TOC convergence instead.
\documentclass[11pt]{article}
\usepackage[T1]{fontenc}
\usepackage{lmodern}
\usepackage[margin=1in]{geometry}

\title{Quarterly Operations Review: Warehouse Efficiency Initiative}
\author{Operations Analytics Team}
\date{Third Quarter, Fiscal Year 2026}

\begin{document}
\maketitle

\section{Headline Figures}

Median order-to-ship time fell by nineteen percent across the network.

Northgate reduced its median cycle time from fourteen hours to ten.

Riverside improved by five percent and remains the slowest facility.

Fairview posted the largest single-quarter gain of the review period.

Union Yard continues to lose afternoon capacity to dock congestion.

Crestline completed its shelving retrofit ahead of the planned schedule.

Ninety-day retention improved at every site that adopted the new model.

Time to full productivity fell from three weeks to ten working days.

Supervisor hours spent on coaching rose by roughly a fifth.

Pick-path reorganization is complete at four of the five facilities.

Two facilities still run the original shelving and routing layout.

Outbound scheduling changes are pending a decision on staggered windows.

A further facility note number 1 is recorded in the appendix.

A further facility note number 2 is recorded in the appendix.

A further facility note number 3 is recorded in the appendix.

\section{Executive Summary}
This report summarizes the findings of the warehouse efficiency initiative
launched at the start of the fiscal year, covering staffing, equipment
utilization, and fulfillment throughput across the five regional
distribution centers. Overall, the initiative has delivered a measurable
improvement in order-to-ship time, though results vary considerably by
site, and several facilities have exceeded their targets while others
continue to struggle with fixture and layout constraints that were only
partially addressed during the mid-year renovation window.

\section{Background}
The warehouse efficiency initiative was chartered in response to a steady
increase in order volume that had begun to outpace the capacity of existing
fulfillment workflows. Waiting times at pick stations had crept upward over
several consecutive quarters, and exit surveys from departing staff
frequently cited workflow friction and unclear task prioritization as
sources of frustration. Management commissioned a cross-functional working
group, drawn from operations, facilities, and workforce planning, to
diagnose the underlying causes and propose a coordinated set of fixes,
and several sites reported the same pattern, and several sites reported
the same pattern, and several sites reported the same pattern.

\section{Throughput Improvements}
Order-to-ship time by facility is summarized in the table below.

\begin{table}[h]
\centering
\caption{Median order-to-ship time by facility (hours).}
\end{table}

\section{Conclusion}
The warehouse efficiency initiative has delivered real, if uneven, gains.

\end{document}
